\begin{figure}[htbp]
\centering
\begin{tikzpicture}[scale=2]
    % Axes
    \draw[->] (-0.5,0) -- (2.5,0) node[right] {$\Real$};
    \draw[->] (0,-0.5) -- (0,2.2) node[above] {$\Imag$};

    % Grid
    \draw[help lines, gray!30] (-0.5,-0.5) grid (2.5,2.2);

    % Unit circle
    \draw[dashed,gray!60,thick] (0,0) circle (1cm);

    % Define coordinates
    % z1 = 0.8 * e^(j*25°) = 0.8(cos25 + j*sin25)
    \coordinate (Z1) at ({0.8*cos(25)},{0.8*sin(25)});
    % z2 = 1.2 * e^(j*40°) = 1.2(cos40 + j*sin40)
    \coordinate (Z2) at ({1.2*cos(40)},{1.2*sin(40)});
    % z1*z2 = 0.96 * e^(j*65°)
    \coordinate (Z1Z2) at ({0.96*cos(65)},{0.96*sin(65)});

    % Draw vectors
    \draw[->,thick,UofSCAtlantic] (0,0) -- (Z1) node[right,UofSCAtlantic,font=\small\bfseries] {$z_1 = r_1\ee^{\jj\theta_1}$};
    \draw[->,thick,UofSCHorseshoe] (0,0) -- (Z2) node[above right,UofSCHorseshoe,font=\small\bfseries] {$z_2 = r_2\ee^{\jj\theta_2}$};
    \draw[->,thick,UofSCGarnet,line width=1.5pt] (0,0) -- (Z1Z2) node[above,UofSCGarnet,font=\small\bfseries] {$z_1 z_2 = r_1 r_2\ee^{\jj(\theta_1+\theta_2)}$};

    % Angle arcs
    % theta1 = 25°
    \draw[->,UofSCAtlantic] (0.35,0) arc (0:25:0.35);
    \node[right,UofSCAtlantic,font=\scriptsize] at ({0.4*cos(12.5)},{0.4*sin(12.5)}) {$\theta_1$};

    % theta2 = 40° (from 0)
    \draw[->,UofSCHorseshoe] (0.5,0) arc (0:40:0.5);
    \node[right,UofSCHorseshoe,font=\scriptsize] at ({0.55*cos(20)},{0.55*sin(20)}) {$\theta_2$};

    % theta1 + theta2 = 65° (from 0)
    \draw[->,UofSCGarnet] (0.65,0) arc (0:65:0.65);
    \node[above right,UofSCGarnet,font=\scriptsize] at ({0.7*cos(32.5)},{0.7*sin(32.5)+0.05}) {$\theta_1+\theta_2$};

    % Magnitude labels with braces
    % r1 magnitude
    \draw[UofSCAtlantic,thick,decorate,decoration={brace,amplitude=3pt,mirror}]
        (0,-0.1) -- ({0.8*cos(25)},-0.1) node[midway,below=3pt,UofSCAtlantic,font=\scriptsize] {$r_1$};

    % r2 magnitude (along the vector direction, offset)
    \node[UofSCHorseshoe,font=\scriptsize,rotate=40] at ({0.6*cos(40)-0.1},{0.6*sin(40)+0.1}) {$r_2$};

    % r1*r2 magnitude
    \node[UofSCGarnet,font=\scriptsize,rotate=65] at ({0.48*cos(65)-0.12},{0.48*sin(65)+0.08}) {$r_1 r_2$};

    % Points at vector tips
    \draw[fill,UofSCAtlantic] (Z1) circle (0.04cm);
    \draw[fill,UofSCHorseshoe] (Z2) circle (0.04cm);
    \draw[fill,UofSCGarnet] (Z1Z2) circle (0.04cm);

    % Origin label
    \node[below left,font=\small] at (0,0) {$0$};

    % Key insight box
    \node[draw,thick,fill=white,font=\small,align=center] at (1.9,1.8) {
        Magnitudes multiply\\
        Angles add
    };
\end{tikzpicture}
\caption{Complex multiplication in polar form. When multiplying two complex numbers $z_1 = r_1\ee^{\jj\theta_1}$ and $z_2 = r_2\ee^{\jj\theta_2}$, the magnitudes multiply ($r_1 r_2$) and the angles add ($\theta_1 + \theta_2$). This geometric interpretation makes polar form ideal for multiplication and division operations.}
\label{fig:complex_multiplication}
\end{figure}
